










\chapter{System diagnostic}

Multiple tests and calibrations need to be done on the experimental system for several purposes. To begin with, we can develop knowledge and skills on operating the setup. In addition, we can study how to verify the authenticity of signals. Meanwhile, during those operations, the properties of the system, such as relaxation times, noise level, flux in gradiometer, can be characterized. In a word, we perform diagnostics to prepare ourselves and the system for axion measurements. 

In this chapter, we start with the calculation of the characteristics of NMR signal. Moreover, we present pulsed-NMR experimental results and analyses, which is the basis for more complicated diagnostics in other sections like sample temperature characterization. Lastly, an overnight measurement gives us an idea of the noise properties as a reference for future dark matter measurement. 

\section{Signal response}\label{sec:Signal_response}

Consider the sample placed at the center of gradiometer where the leading field is perpendicular to the orientation of gradiometer coils (one such example is illustrated in Fig.\,\ref{fig:Gradiometer_Sample_Dimensions}(a)). 
The sample is able to generate time-dependent magnetic flux after it is excited by rf pulse or due to spin projection noise (SPN). The response in the reception system is analyzed in this section. 

\subsection{Flux in gradiometer}

We can roughly estimate the flux in a loop by 
\begin{align}
    \Phi_{loop} = B A_{loop}\,,
\end{align}
where $\Phi_{loop}$ is the flux, $B$ is the magnetic field (perpendicular to loop plane) and $A_{loop}$ is the area of the loop. However, the magnetic field generated by the sample is not a homogeneous magnetic field to the gradiometer near it since the sample holder occupies a large volume and the gradiometer is not far from the sample compared to the dimensions of sample holder. Besides, the coils of gradiometer are wrapped on a cylindrical plane. To account for the two points mentioned above, we introduce more complicated calculation on the gradiometer flux in this section. 

\subsubsection{Analytical analysis for small sample}\label{sec:Signal_response:flux:small}

We begin with one assumption to analyze the flux that the spins generate in the gradiometer: the dimensions of the sample is relatively small in comparison to those of the gradiometer. In other words, we assume that every spin stays at same position and generates same amount of flux in gradiometer. Here we define z-axis as the axial orientation so that x-y plane is the transverse plane (see Fig.\,\ref{fig:Gradiometer_Sample_Dimensions}(a)). 
Consider a sample with transverse magnetization $\mathbf{M}_{xy}$ and magnetic dipole moment $\mathbf{m}_{xy} = M_{xy} V$, where $V$ is the volume of the sample. Similarly, the axial magnetization and magnetic dipole moment are $\mathbf{M}_{z}$ and $\mathbf{m}_{z} = M_{z} V$. 
We can write it down in vector form:
\begin{equation}\label{eq:smallsample:m_s}
    \mathbf{m}_s= V(M_{xy}\cos{\xi}, M_{xy}\sin{\xi}, M_{z})\,,
\end{equation}
where $\xi$ is the azimuth of the vector (angle on the transverse plane). 
The magnetic field that the sample generates $\mathbf{B_s}$ at $\mathbf{r}=(R\cos\theta,R\sin\theta,z)$ is determined by:
\begin{equation}\label{eq:smallsample:B_s}
    \mathbf{B_s}(\mathbf{r})=\frac{\mu_0}{4\pi} \frac{3(\mathbf{m}_s\cdot\mathbf{r})\mathbf{r} - |\mathbf{r}|^2 \mathbf{m}_s}{|\mathbf{r}|^5}\,,
\end{equation}
% Saddle Coils for Uniform Static Magnetic Field Generation in NMR Experiments
% https://onlinelibrary.wiley.com/doi/pdf/10.1002/cmr.b.20057
where $\mu_0$ is vacuum permeability. The corresponding flux in gradiometer is:
\begin{equation}\label{eq:smallsample:Phi}
    \Phi = \int_{S} \mathbf{B_s}\cdot d\mathbf{S}\,,
\end{equation}
where the derivative of vector area of the saddle coil $d\mathbf{S}$ at $\mathbf{r}$ is
\begin{align}
    d\mathbf{S}=dS\mathbf{n}=R_c d\theta dz (\cos\theta, \sin \theta, 0)\,.\label{eq:Coil_dS}
\end{align}
Here $\mathbf{n}$ is the normal vector of the coil surface at $\mathbf{r}$. $R_c$ is the radius of the gradiometer. $\theta$ is the angle of elevation for $\mathbf{r}$. 
Note that the orientation magnetic field is perpendicular to gradiometer surface, thus $\mathbf{m}_{z}$ does not generate flux in gradiometer.  
The illustration of it is shown in Fig.\,\ref{fig:Gradiometer_Sample_Dimensions}(b). 

\begin{figure}[htbp]
    \includegraphics[width=\textwidth]{figures/03_Diagnostic/Gradiometer_Sample_CoordinateFrame.png}
    \caption{(a) Layout and dimensions for gradiometer and sample. (b) (c) Schematic for integral of flux in a coil. A small sample is placed at the center of gradiometer in (b). While in (c), it is placed at an arbitrary position. }
    \label{fig:Gradiometer_Sample_Dimensions}
\end{figure}

Assembly Eqs.\,\eqref{eq:smallsample:m_s}, \eqref{eq:smallsample:B_s} and \eqref{eq:smallsample:Phi}, so that for one half of a saddle coil (e.g., the green coil in Fig.\,\ref{fig:Gradiometer_Sample_Dimensions}(b)) we obtain
\begin{equation}
    \Phi_{1half} (\theta_0,\theta_1, R, z_0, z_1) =
    \frac{\mu_0 M_{xy} V }{4\pi} 
    \sin(\theta-\xi)\bigg|^{\theta_1}_{\theta_0} \, 
    \frac{z (2 R^2 + z^2)}{R(R^2 + z^2)^{3/2}} \bigg|^{z_1}_{z_0}
    \,.\label{eq:Phi1half}
\end{equation}
Looking at the projection of one piece of saddle coil pair onto the transverse plane, $\theta_0$ and $\theta_1$ are the angles of azimuth of two ends of the arc (see top view in Fig.\,\ref{fig:Gradiometer_Sample_Dimensions}(a)), therefore $\theta_1-\theta_0$ is the angle of the arc. $R$ is the radius of it. $z_0$ and $z_1$ are the z-coordinates of the lower and upper arcs in the saddle coil. 

Based on Eq.\,\eqref{eq:Phi1half}, the flux in a saddle coil pair can be drawn as:
\begin{align}
	\Phi_{sd} (R, z_0, z_1) &= 
	\Phi_{1half} (\theta_0=\ang{30} ,\theta_1=\ang{150}) + \Phi_{1half} (\theta_0=\ang{-30} ,\theta_1=\ang{-150})\nonumber\\
	&=\frac{\sqrt{3}}{2\pi}\mu_0 M_{xy} V  \sin (\xi) \,
	\frac{z (2 R^2 + z^2)}{R(R^2 + z^2)^{3/2}} \bigg|^{z_1}_{z_0}%\\
	%&=\mathrm{g}_V \mu_0 M_{xy} V  \sin (\xi)\label{eq:smallsample:gradflux}
	\,.
\end{align}
% Parameters $R$, $z_0$ and $z_1$ are omitted in $\Phi_{1half}$ expressions here only for brevity. Details on gradiometer were written in \S\ref{sec:Signal_receptio_and_processing:Gradiometer_and_sample_holder}. 

The values for dimensions are listed in Fig.\,\ref{fig:Gradiometer_Sample_Dimensions}(a). 
From the top to the bottom, the flux in three saddle coils are:
\begin{align}
    \Phi_{sd1} &= (\SI{0.33}{\meter^{-1}}) \mu_0 M_{xy} V \sin (\xi) \\
    \Phi_{sd2} &=(\SI{43.64}{\meter^{-1}}) \mu_0 M_{xy} V \sin (\xi) \\
    \Phi_{sd3} &=(\SI{-0.32}{\meter^{-1}}) \mu_0 M_{xy} V \sin (\xi)\,. 
\end{align}
From the values in these three equations, we can learn that the fluxes in side coils are much less than the one in the central coil. Meanwhile, total flux in the gradiometer is 
\begin{align}
    \Phi_{grad}&=-\Phi_{sd1}+\Phi_{sd2}-\Phi_{sd3} \nonumber \\
                &=(\SI{43.63}{\meter^{-1}}) \mu_0 M_{xy} V\sin (\xi)\,. \label{eq:Phigrad_smallsample}
\end{align}
Here the plus and minus symbols indicate the orientation of the saddle coils as discussed in \S\,\ref{sec:Signal_receptio_and_processing:Gradiometer_and_sample_holder}. 

We can further summarize the characteristics of $\Phi_{grad}$ into:
\begin{align}
    \Phi_{grad}&=\mathrm{g}_V \mu_0 M_{xy} V \sin (\xi)\,. \label{eq:Phigrad_gV}
\end{align}
In this equation, $\mathrm{g}_V$ is called geometry coupling constant which is determined by the layout and dimensions of gradiometer and sample. 

\subsubsection{Numerical analysis for large sample}

From here we take the volume of sample into consideration. In other words, the nuclear spins now have spatial distribution. 
We first follow a similar pattern as in \S\,\ref{sec:Signal_response:flux:small}. Consider a tiny fraction of sample placed at 
\begin{align}
    \mathbf{r}_s =(x_s , y_s, z_s)\,.
    % =(R_s\cos\theta_s, R_s\sin\theta_s,z_s)
\end{align}
% Here $R_s = \sqrt{x_s^2 + y_s^2}$, $\theta_s = \arctan(y_s/x_s)$. $= R_s dR_s d\theta_s dz_s$
The volume of this fraction is $dV = dx_s dy_s dz_s$. We denote sample magnetization as $\mathbf{M}_s$, which is density of sample magnetic moment $\mathbf{m}_s$: 
\begin{equation}\label{eq:largesample:m_s}
    \dfrac{d\mathbf{m}_s}{dV} = \dfrac{\mathbf{m}_s}{V} = \mathbf{M}_s = (M_{xy}\cos{\xi}, M_{xy}\sin{\xi}, M_{z}) \,.
    % =(\rho_{m_{xy}}\cos{\xi}, \rho_{m_{xy}}\sin{\xi},  \rho_{m_z}) R_s dR_s d\theta_s dz_s,
\end{equation}
Still, $M_{xy}$ and $M_{z}$ stand for the transverse and axial magnetization respectively. $\xi$ is the azimuth of the vector. 
The magnetic field it generates at $\mathbf{r}$ is given by
\begin{equation}\label{eq:largesample:B_s}
    \mathbf{B_s}(\mathbf{r})=\int_{V} \frac{\mu_0}{4\pi} \frac{3(\mathbf{M}_s\cdot \Delta \mathbf{r})\Delta \mathbf{r} - |\Delta \mathbf{r}|^2 \mathbf{M}_s}{|\Delta \mathbf{r}|^5} dV\,,
\end{equation}
where $\Delta \mathbf{r}$ stands for $\mathbf{r} - \mathbf{r}_s$. 

The corresponding flux in a coil is
\begin{equation}
    \Phi = \int \mathbf{B_s}(\mathbf{r})\cdot d\mathbf{S}(\mathbf{r})\,, 
\end{equation}
where $d\mathbf{S} $ is consistent with the definition in Eq.\,\eqref{eq:Coil_dS}. 

With a given sample holder, flux in gradiometer can still be given in the form of Eq.\,\eqref{eq:Phigrad_gV}. Numerical integration can be done to compute the value of $\mathrm{g}_V$. 
Furthermore, we can examine the numerical integration in such a way: Set the sample holder dimensions to values much smaller than those of gradiometer, compute $\mathrm{g}_V$, and compare it to the result under small-sample assumption. They should agree within the error due to numerical integration. Otherwise we know there are problems in the calculation. 

The values for $\mathrm{g}_V$s are summarized in Tab.\,\ref{tab:gV}. When we use numerical method for computing the $\mathrm{g}_V$ of a tiny sample ($\sim \SI{3}{\nm^{3}}$), the result agrees with the analytical solution in 1\%. This error can be attributed to the low accuracy goal setting in the calculation via Mathematica. Moreover, we find in the comparison between small-sample and large-sample results that the spatial distribution of sample decreases the magnitude of $\mathrm{g}_V$. 
% see casper-gradient-code\Supplementary\Gradiometer Flux\flux in gradiometer.nb
\begin{table}[ht]
    \centering
    \begin{tabular}{cccc}
    	\hline
    	\multirow{2}*{Sample} & \multirow{2}*{Infinitely small} & Cylindrical & Cylindrical \\
    	 & & {\footnotesize $R=\SI{1}{\nm}, H=\SI{1}{\nm}$} & {\footnotesize $R=\SI{4}{\mm}, H=\SI{22.53}{\mm}$} \\
    	\midrule
    	Method & Analytical solution & \multicolumn{2}{c}{Numerical integration}  \\
    	\midrule
    	$\mathrm{g}_V\,/\,\SI{}{\meter^{-1}}$ & 43.6323 & 43.54 & 37.67  \\
    	\bottomrule
    \end{tabular}
    \caption{Geometry coupling constant $\mathrm{g}_V$ obtained under different sample size. $R$ and $H$ are radius and height of a cylindrical sample. The third sample ($R=\SI{4}{\mm}, H=\SI{22.53}{\mm}$) is the sample in use. The center of the sample is at the center of the gradiometer. }
    \label{tab:gV}
\end{table}

\subsection{Flux in SQUID}

The coupling from gradiometer to SQUID is discussed in \S\,\ref{sec:Signal_receptio_and_processing:SQUID}. Here we take advantage of Eqs.\,\eqref{eq:Phiin_Phigrad} and \eqref{eq:Phigrad_gV}, and the input flux in SQUID loop is
\begin{align}
	\Phi_{in} = \mathrm{g}_{in} \mathrm{g}_V \mu_0 M_{xy} V \sin\xi / L\,, \label{eq:Phiin_gV}
\end{align}
where $L = L_{grad} + L_{in}$ is the loop inductance. 

The magnitudes of $M_{xy}$ depends on the experimental scheme. Here we compute input flux for pulsed-NMR and SPN. 

\subsubsection{Pulsed-NMR flux with thermally-polarized sample}

Consider an ideal $\pi/2$ pulse which tips the magnetization $\mathrm{M} = M_0 \mathrm{e}_z$ by \SI{90}{\degree}. Take Eqs.\,\eqref{eq:M0_allspins} and \eqref{eq:MxyM0theta}, then we can have
\begin{align}
    M_{xy} = M_0 = \mu p n\,,\label{eq:Mxy_90deg_mu_p_n}
\end{align}
where $M_0$ is the magnetization magnitude at thermal equilibrium, $\mu$ is the nuclear magnetic moment, $p$ is polarization, and $n$ is the spin density. 

The magnetization vector precesses along the leading field $\mathbf{B} = \mathrm{B} \mathbf{e}_z$, therefore the azimuth $\xi$ is time-dependent:
\begin{align}
    \xi = \omega t + \phi_0\,,
\end{align}
where angular frequency $\omega = \gamma \mathrm{B}$, $t$ is time and $\phi_0$ is the initial phase. The amplitude of $\Phi_{in}$ can be inferred from Eqs.\,\eqref{eq:Phiin_gV} and \eqref{eq:Mxy_90deg_mu_p_n}:
\begin{align}
	\Phi_{\pi/2} = \mathrm{g}_{in} \mathrm{g}_V \mu_0 \mu p n V / L\,, \label{eq:Phiin_gV_pNMR}
\end{align}

To better demonstrate the details on calculation, we list the information of parameters;
\begin{itemize}
    \item Geometry coupling constant $\mathrm{g}_V$. For cylindrical sample ($R=\SI{4}{\mm}, H=\SI{22.53}{\mm}$) the value of it is \SI{37.67}{\meter^{-1}}
    \item Vacuum permeability $\mu_0$. The value of it is \SI{1.256 637 062 12(19)e-6}{\henry.\metre^{-1}}\cite{FUNDAMENTALCONSTANTS}. 
    \item Nucleus magnetic moment $\mu$. Here the nucleus is proton. The value of it is \SI{1.410 606 797 36(60)e-26}{\joule.\tesla^{-1}}\cite{FUNDAMENTALCONSTANTS}
    \item Thermal polarization $p$. We take temperature of \SI{-90}{\celsius}, leading field of \SI{317}{\gauss}, proton gyromagnetic ratio \SI{42.577 482 5}{\MHz. \tesla^{-1}}\cite{rumble2019crc}, and $p$ is computed to be \SI{1.78e-7}{} by Eq.\,\eqref{eq:thermalp_approx}. 
    \item Spin density $n$ for methanol is \SI{112.87}{\mmol.\cm^{-3}} (taken from Tab.\,\ref{tab:samplesproperity}). 
    \item Cylindrical sample volume $V = \pi R^2 H = \SI{1.132}{\cm^{3}}$. 
    \item The coupling strength between SQUID flux and gradiometer flux is given by $\mathrm{g}_{in} / L$. For SQUID sensor C6L1W and C73L1, this values are \SI{6.88e-5}{} and \SI{7.14e-5}{} respectively (taken from Tab.\,\ref{tab:SQUIDsensors}). 
\end{itemize}

With the values listed above, we can compute pulsed-NMR flux in the SQUID loop for two sensors (omitting the precession term $\sin\xi$):
\sisetup{exponent-mode = scientific}
\begin{align}
	\Phi^{C6L1W}_{\pi /2} &=\SI{0.000264992}{}\,\Phi_0\\
	\Phi^{C73L1}_{\pi /2} &=\SI{0.000272598}{}\,\Phi_0\,. 
\end{align}
\sisetup{exponent-mode=input}

The positions of sample during different measurements may not be consistent. We need a method to diagnose the position of sample center with regard to the gradiometer center. The flux is one parameter to infer the position. When the sample is moved from the center to the gaps between saddle coils, the flux decreases. It increases again when it is further moved to the center of a side coil. 


Besides experimental measurement, such behavior can be estimated with numerical integration method we developed in this section. The results are shown in Fig.\,\ref{fig:flux-samplez}

\begin{figure}[ht]
	\centering
	\includegraphics[width=\textwidth]{figures/03_Diagnostic/Flux_in_Gradiometer_20220915_3.png}
	\caption{Flux in SQUID with respect to the position of sample z-coordinates (axial direction). (a) Experimental data before correcting the x-axis [mm]. (b) Blue curve: numerical integration result. The curve is not smooth due to computation errors. Orange scatters in (b): shifted experimental results.}
	\label{fig:flux-samplez}
\end{figure}

The sample was lifted along z-axis to obtain data at 14 different sample positions. The distance of its starting point from gradiometer center along z-axis was unknown before the measurement. 
We can observe an offset of $\sim \SI{1.8e-4}{}\,\Phi_0$ in $\Phi_{in}$ from experimental results, which is attributed to the noise power in the spectrum. If we cancel this offset, and shift the x-axis of measurements by \SI{-15.15}{\mm}, we can find the measurements roughly agree with numerical method results as shown in Fig.\,\ref{fig:flux-samplez}(b), which indicates the starting point of the measurement is \SI{15.15}{\mm} lower than the gradiometer center. In this way we can optimize sample position for large signal. 

In addition to the flux in the SQUID, we are also interested in the signal amplitude on spectra as discussed in \S\,\ref{sec:NMR:Pulsed-NMR}. Here we choose \SI{30}{\Hz} linewidth for the signal and the amplitudes can be estimated through Eq.\,\eqref{eq:Lorz_nuc}:
\sisetup{exponent-mode = scientific}
\begin{align}
	S^{C6L1W}_{\pi /2} &=\SI{1.49013e-9}{}\,\Phi_0^{2}/\SI{}{\Hz}\\
	S^{C73L1}_{\pi /2} &=\SI{1.5769e-9}{}\,\Phi_0^{2}/\SI{}{\Hz}\,. 
\end{align}
\sisetup{exponent-mode=input}

\subsubsection{Spin-projection noise}\label{sec:SR:FluxinSQUID:SPN}

According to  Eq.\,\eqref{eq:SPN_M} we take $M_{xy} =\sqrt{N/2\pi} \mu/V $ for Eq.\,\eqref{eq:Phiin_gV} to estimate flux generated by SPN:
\begin{align}
	\Phi_{SPN} = \mathrm{g}_{in} \mathrm{g}_V \mu_0 \sqrt{N/(2\pi)} \mu \sin\xi / L\,. \label{eq:Phiin_gV_SPN}
\end{align}
Now we are able to compute its magnitude for our sensors:
\begin{align}
	\Phi^{C6L1W}_{SPN} =\SI{2.30e-9}{}\,\Phi_0\\
	\Phi^{C73L1}_{SPN} =\SI{2.37e-9}{}\,\Phi_0\,. 
\end{align}
Furthermore, the signal amplitudes for SPN for sensors are:
\begin{align}
	S^{C6L1W}_{SPN} =\SI{1.12e-19}{}\,\Phi_0^{2}/\SI{}{\Hz}\\
	S^{C73L1}_{SPN} =\SI{1.19e-19}{}\,\Phi_0^{2}/\SI{}{\Hz}\,. 
\end{align}
SPN should be taken into account when the excitation field for NMR is relatively small. We come back to it during the system noise analysis in \S\,\ref{sec:Noise_analysis}. 

\section{Pulsed-NMR spectroscopy}

In \S\,\ref{sec:NMR} we mentioned relaxation of NMR signal in time domain and the related broadening in frequency domain. Besides, chemical environment and spin-spin coupling inside sample molecule could shift nuclear resonant frequency. In this section we take methanol as an example to demonstrate basic pulsed-NMR spectroscopy analysis. 

Regular methanol consisting of \ch{^{1}H}, \ch{^{12}C} and \ch{^{16}O}. Only \ch{^{1}H} nucleus among them has nonzero nuclear spin, with which NMR signal can be generated. We first look into the spectrum of it. In addition to that, isotope-enriched methanols are adopted for special purposes. For instance, deuterium (\ch{^{2}H}) has nuclear spin of 1. The NMR frequency of it differs from that of proton. Therefore, \ch{^{2}H}-enriched methanol helps us to study NMR properties with individual proton in the molecule. \ch{^{13}C}-enriched methanol spectra are analyzed in details here, which assist us in magnetic field measurement and signal verification. 

\subsection{Regular methanol}

There are two types of protons in methanol: ones in methyl (\ch{CH3}) and the one in hydroxyl (\ch{OH}). The chemical shift between these two groups is \SI{1.5}{\ppm} (illustrated in Fig.\,\ref{fig:PulsedNMR_Methanol}(a)). To resolve this shift on spectrum, the linewidths of two signal peaks should not be greater than \SI{1.5}{\ppm}. 


We acquired a few PSD spectra under same setting for main coil and different settings for shim coils (see Fig.\,\ref{fig:PulsedNMR_Methanol}(b)). The resonant frequency is approximately \SI{1.35}{\MHz}. The \SI{1.5}{\ppm} chemical shift translates into \SI{2}{\Hz}, which is much smaller than the linewidths here. Therefore we know such broadening cannot be attributed to chemical shift. Other reasons for the broadening are analyzed as follows. 
\begin{figure}[htbp]
    \includegraphics[width=\textwidth]{figures/03_Diagnostic/12CMethanol_CS_and_Inhomo_20220912.png}
    \caption{(a) Simulated methanol NMR spectrum with chemical shift. (b) Methanol spectra under different shim field settings. }
    \label{fig:PulsedNMR_Methanol}
\end{figure}

$T_2$ time for liquid methanol is approximately \SI{2}{\second} at room temperature, corresponding to a linewidth of $\SI{0.2}{\Hz}$. When the resonant frequency is greater than \SI{0.2}{\MHz}, the chemical shift is larger than the linewidth due to $T_2$ relaxation, which makes it easier to resolve the signal peaks on spectrum. This is one of the reasons why chemists prefer high-field magnet. However, due to imperfect insulation inside cryostat, the sample is approximately \SI{-90}{\celsius}, well below room temperature. Spin echo experiment was performed to measure $T_2$. Analysis on this part is included in \S\,\ref{sec:SpinEchoExperiment}. The $T_2$ time was measured to be greater than \SI{20}{\ms}, indicating that linewidth $1/(\pi T_2)$ should be no greater than \SI{16}{\Hz}, which is still smaller than spectral linewidth ($>\SI{22}{\Hz}$). This implies that other mechanisms contributed to the broadening. 

Inhomogeneous broadening, due to inhomogeneous leading field (accounted in relaxation time $T_{\delta}$), is another obstacle for achieving high resolution. Those different lineshapes of spectra demonstrates that we are able to manipulate $T_{\delta}$ with shim coils. By means of, for example, nonlinear fitting, we can characterize signal signature and provide information for shimming algorithm. 

\subsection{\ch{^{13}C}-enriched methanol}

Considering the common practice in chemistry society, we adopt normalized linewidth here to analyze spectra. TMS is widely used as the reference sample since there is no chemical shift. Note that normalized linewidth increases from right to the left. 

Because of the coupling between nuclei of \ch{^{13}C} and \ch{^{1}H}, a doublet state can be composited. When the nuclear spins of \ch{^{13}C} and \ch{^{1}H} are parallel, since both of their magnetic moment are positive, they are at lower energy state and vice versa. Therefore, for proton NMR, the NMR signal from methyl splits into two. The shift with regard to the \ch{^{12}C^{1}H3} NMR frequency are $\pm\SI{70}{\Hz}$. Depending on the reference frequency, this splitting is transformed into different normalized linewidth values. We plot a few \ch{^{13}C}-methanol spectra under different leading field in Fig.\,\ref{fig:C13Methanol}. When at low field, i.e., when when splitting of \SI{140}{\Hz} is more significant compared to the chemical shift, two \ch{^{13}CH3} proton NMR peaks are located at two sides of \ch{OH}.
On the other hand, when the field is so high that chemical shift is larger than half of the splitting (\SI{70}{\Hz}), the \ch{OH} signal peak stays on one side of \ch{^{13}CH3} peaks. 
% We can find such examples in Fig.\,\ref{fig:C13Methanol}(a) and (b). 
% Such examples are \SI{0.5}{\tesla} simulated spectrum and Fig.\,\ref{fig:C13Methanol}(b) obtained at \SI{316}{\gauss}. 
% $\ch{^{13}C}(\uparrow)-\ch{^{1}H}()$, $\ch{^{13}C}(\uparrow)-\ch{^{1}H}(\downarrow$)
\begin{figure}[htp]
    \includegraphics[width=\textwidth]{figures/03_Diagnostic/C13Methanol_Jcoupling_20220912_5.png}
    \caption{\ch{^{13}C}-labeled methanol spectra under different leading field. All black arrows are indicating the \SI{140}{\Hz} splitting due to spin-spin coupling. Simulated spectra are shown in (a). Spectra obtained at \SI{316}{\gauss} and \SI{1.4}{\tesla} are shown in (b) and (c) separately. }
    \label{fig:C13Methanol}
\end{figure}

During experiment where the leading field is unknown and unshimmed, signal-noise-ratio (SNR) can be low. Besides, noise spikes could resemble regular methanol signal in terms of amplitude and linewidth, especially when the $T_2^*$ time is so short that we cannot resolve the chemical shift. 
Since \ch{^{13}C}-labeled methanol spectrum has more characteristics which are not possessed by noises, it is easier to decrease the influence of noise. 
The \SI{140}{\Hz} splitting creates a trident-like lineshape, which is a unique signature for signal verification. The amplitudes of these three peaks can be computed based on the number of protons in methyl and hydroxyl. These are evidences identifying the NMR signal. 


\section{Alignment to magnetic center}

The geometric (also magnetic) centers of main coil and shim coils are overlapping. The field homogeneity becomes worse at positions far from the center. One such example of Helmholtz coil field is illustrated in Fig.\,\ref{fig:FOGS}(a) to (c). 

Superconducting magnet is placed deep down in cryogenic vessel, as a consequence of which the sample insert is relatively long (approximately \SI{2}{m} for our experiment) to reach the magnetic center. Even with a carefully-designed setup, aligning sample with the magnetic center within \SI{1}{\mm} accuracy still requires extra efforts due to the length of the insert. Fine adjustment to sample position is crucial. 

Here, a method inspired by magnetic resonance imaging (MRI) is introduced to figure out the relative sample position with regard to the magnetic center. It is named first order gradient scrutinizing (FOGS) since it involves first-order shim channels (X, Y and Z1). These channels can produce a $\mathbf{B}_z$ field, magnitudes of which have linear dependence on  spatial coordinates $x$, $y$ and $z$ respectively. An example of channel Z1 is displayed in Fig.\,\ref{fig:FOGS}(d) and (e). More details of these channels are included in Tab.\,\ref{tab:shimproperity}. 

In the scheme of the method, we first obtain a well-shimmed NMR signal. The next step, is to deliberately apply additional current in shim channel X, Y or Z1, corresponding to alignment on x, y or z axis. These channels generate zero field at the center. Therefore, when sample is placed at the center of the related axis, the central frequency of NMR cannot be shifted when we increase the current in the shim channel. Nevertheless, when the sample is not placed at the center, the magnitude of leading field increases with the current in the shim channel, resulting in the shift of NMR central frequency. 
Another advantage of this method is that, the NMR signal is shifted towards either higher or lower frequency depending on the sample position offset. This provides information for adjusting sample position. The scheme is drawn in Fig.\,\ref{fig:FOGS}(f) to (h) as well. 
\begin{figure}[htbp]
  \centering
    \includegraphics[width=\textwidth]{figures/03_Diagnostic/Z1sweep_Theo_20220914.png}
  \caption{(a) Schematic for Helmholtz (green) and anti-Helmholtz (red) coil. The fields they generate are plotted in (b) to (e). Sample (gray sphere) placed with different z-offset are illustrated in (f), (g) and (h). The magnetic fields have linear gradient due to shim channels. With nonzero z-offsets, signal frequency is shifted towards either larger or smaller values when the current in the shim channel is increased. }
  \label{fig:FOGS}
  % Biot-Savart Law assuming 
\end{figure}

Moreover, we can quantify the relationship between the frequency shift and sample position offset. The frequency shift $\Delta \nu$ is proportional to the change in leading field $\Delta \mathrm{B}$:
\begin{align}
    \Delta \nu = \dfrac{\gamma}{2 \pi} \Delta \mathrm{B} \,,\label{eq:Deltanu_DeltaB}
\end{align}
where $\gamma$ is the gyromagnetic ratio. Here $\Delta \mathrm{B}$ is produced by the current variation $\Delta I$ in the shim channel. We have
\begin{align}
    \Delta \mathrm{B} = \mathrm{g}_{m} \Delta I (r - r_c)\,,\label{eq:DeltaB_DeltaI_r},
\end{align}
where $\mathrm{g}_{m}$ is the factor to describe shim field configuration, and $r$ is the spatial coordinate, while $r_c$ is the magnetic center coordinate. We can establish the formula with Eqs.\,\eqref{eq:Deltanu_DeltaB} and \eqref{eq:DeltaB_DeltaI_r} to estimate sample position:
\begin{align}
    \Delta \nu = \dfrac{\mathrm{g}_{m} \gamma}{2 \pi} \Delta I (r - r_c)\,.\label{eq:Deltanu_DeltaI_r}
\end{align}
This equation shows that the shift in NMR frequency is proportional to the sample coordinate and current variation. Measurements were done to demonstrate this method in practice. The air spring on VTI lifted up the gradiometer and sample together to a few different positions. At each position, different currents were applied in Z1 channel. We obtained spectra under such settings. Frequency shifts can be observed in Fig.\,\ref{fig:Z1sweep_Exper}(a). We plot the shift of central frequency with regard to Z1 current in Fig.\,\ref{fig:Z1sweep_Exper}(b). The slopes of the curves can help us to determine the value of $\mathrm{g}_{m}\Delta I (r + r_c)/(2\pi)$ in Eq.\,\eqref{eq:Deltanu_DeltaI_r}. 
However, we do not need to compute it if the purpose is to align the sample to the magnetic center. If we look at $\Delta \nu$ obtained under same $\Delta I$, a linear fitting on $\Delta \nu$ and $r$ can be done. The fit curve crosses the x-axis at $r = r_c$. In this way we can deduce the coordinates of magnetic center. 
Related experimental results are shown in Fig.\,\ref{fig:Z1sweep_Exper}(c). 
\begin{figure}[htbp]
    \centering
    \includegraphics[width=\textwidth]{figures/03_Diagnostic/Z1sweep_Exper_20220916_3.png} 
    \caption{\ch{^{13}C}-enriched methanol spectra obtained different Z1 current offsets are plotted in (a) to (d). The colors indicate the offset in gradiometer and sample position in z direction, which are illustrated in (e). (f) Frequency shift due to Z1 current offset.  (g) Here difference of signal central frequency at $\Delta I = \SI{.75}{\ampere}$ and $= \SI{-.25}{\ampere}$ is taken as the frequency shift. The gray line is the linear fitting based on measurements at air spring z-offset = \SI{0}{\mm}, \SI{18}{\mm} and \SI{29}{\mm}. Later the air spring was adjusted to \SI{14.5}{\mm}, which was predicted to be the magnetic center by the fitting result. }
    \label{fig:Z1sweep_Exper}
\end{figure}

\section{Spin echo experiment}\label{sec:SpinEchoExperiment}

When we need large signal amplitude, we prefer long relaxation time. $T_2^*$ can be estimated from pulsed-NMR signal linewidth. However, we cannot distinguish $T_2$ and $T_\delta$ by this method. Since $T_2^*$ is the inverse of the sum of 
$1/T_2$ and $1/T_\delta$ (see Eq.\,\eqref{eq:T2star_T2_Tdelta}), either $T_2$ or $T_\delta$ can be the bottleneck in increasing $T_2^*$. 

We performed spin echo experiment to figure out the which one of them two is limiting the signal linewidth. The measurement was repeated under different temperatures, which was characterized by a thermocouple thermometer inserted in the sample holder. A nonlinear fitting to function $f = A e^{-t/T_2}$ was done to obtain the $T_2$ time. Fitting results and the dependence on temperature are plotted in Fig.\,\ref{fig:SpinEcho_experiment}. 
\begin{figure}[htbp]
    \centering
    \includegraphics[width=\textwidth]{figures/03_Diagnostic/T2_temperature_Analysis_20220916.png}
    \caption{(a) Curve fitting of echo amplitude and time. (b) Gray dash line and scatters on it: $T_2$ times computed from curve fitting. Blue dash line and scatters on it: The linewidth due to $T_2$ relaxation. }
    \label{fig:SpinEcho_experiment}
\end{figure}

We can figure out that $T_2$ has temperature dependence. For our experimental environment (\SI{-60}{\celsius} to \SI{-20}{\celsius}), $T_2$ ranges from \SI{20}{} to \SI{40}{\ms}, corresponding to a linewidth of $< \SI{16}{\Hz}$. If the signal linewidth is much larger than this value under similar temperature, we can infer that inhomogeneous broadening plays an more important role in relaxation. 

One explanation on the temperature dependence is the microscopic magnetic field homogeneity. In addition to macroscopic magnetic field, spins experience interactions with neighboring spins. Such interaction between spins forms the the microscopic magnetic field. The spins in liquid or gas phase are moving so fast that the average effect is subtle. When sample is at or close to solid state, the inhomogeneity of microscopic magnetic field becomes significant. This results in shorter $T_2$ time at lower temperature. 
Methanol NMR spectra close to melting point (\SI{175.61}{\kelvin}) and related discussion are included in \S\,\ref{sec:stability} and Fig.\,\ref{fig:FreezingTest}. 

\section{System stability characterization}\label{sec:stability}

Axion measurement can last for hours, which requires control and monitoring over sample temperature and leading field magnitude. They could influence the measurement in these aspects:
\begin{enumerate}
    \item Polarization. Thermal polarization is inversely proportional to temperature. Thus, the fluctuation in temperature is related to the fluctuation in expected signal amplitude. 
    \item Chemical shift. Temperature influences the chemical environment and consequently, the chemical shift. The NMR signature could be different due to it. 
    \item $T_2$ relaxation. As mentioned in \S\,\ref{sec:SpinEchoExperiment}, the motion of neighboring spins produces the microscopic magnetic field for spins. The homogeneity of such field becomes worse as spins slow down at lower temperature. Inhomogeneity leads to smaller $T_2$ time and larger signal linewidth. 
    \item Drift in leading field magnitude. If the magnitude is not steady overtime, the resonant frequency changes accordingly. 
\end{enumerate}

Here we still take methanol sample as the example. The first issue on polarization variation during the measurement is easy to solve: we apply pulse at certain time interval to monitor the signal properties including polarization. Meanwhile, nitrogen gas can be blown to maintain the temperature of sample in the cryostat. If there is no action in temperature control, due to the thermal contact and black body radiation, the sample cools down to \SI{-110}{\celsius}, which is below the melting point of methanol. Nevertheless, we performed a 3-hr measurement without nitrogen flow to study polarization of sample continuously cooling down. Besides, we should observe broadening when methanol is freezing due to $T_2$ relaxation. 

Concerning the chemical shift, its temperature dependence is plotted in Fig.\,\ref{fig:FreezingTest}(a) based on \cite{VanGeet_methanol, Geet1970CalibrationOM}. 
The pulsed-NMR linewidth is approximately \SI{40}{\Hz} at \SI{1.349530}{\MHz}, and the corresponding normalized linewidth of NMR is \SI{30}{\ppm}, which is 10 times larger than chemical shift of liquid methanol. Hence we do not have to consider chemical shift here. 

Moreover, drift in superconducting magnet is not supposed to happen within a few days unless scenarios like losing superconductivity. We can monitor it with pulsed-NMR diagnostic as well. 

\begin{figure}[bhtp]
    \centering
    \includegraphics[width=\textwidth]{figures/03_Diagnostic/FreezingTest_Results_20220917.png}
    \caption{(a) Chemical shift of regular methanol under different temperatures. Reproduced from \cite{VanGeet_methanol, Geet1970CalibrationOM}. (b) Pulsed-NMR spectra over time. Here $\nu_c$ is \SI{1.349530}{\MHz}. (c), (d), (e) and (f) are flux in SQUID, signal amplitude, FWHM (linewidth) and central frequency over time respectively. }
    \label{fig:FreezingTest}
\end{figure}

The measurement took 5 hours, results of which are presented in Fig.\,\ref{fig:FreezingTest}(b) to (f). 
We can infer from the signal amplitude that the sample got frozen at \SI{200}{\minute} during measurement. 
The flux in SQUID is computed by integration over signal peak. After the sample was frozen, the flux remains at approximately $\SI{1.26e-4}{}\,\Phi_0$ due to the noise power. If we cancel this offset, the flux in SQUID contributed by nuclear spins ranges from $\SI{1.48}{}$ to $\SI{2.62}{}\,\Phi_0$. The ratio of the maximal flux to the minimal here is 1.77. We can compare this to the ratio of temperatures. The maximal thermal polarization is obtained at melting point \SI{175.61}{\kelvin}, and we take room temperature as \SI{293.15}{\kelvin}. So that the ratio of polarization at melting point to that at room temperature is computed to be 1.67. This ratio is close to the one we obtain by measuring the flux. Therefore the signal amplitude increase can be attributed to larger polarization at lower temperature. 

Linewidth and central frequency of signals are estimated by curve fitting to Lorentzian function. 
Due to the low SNR, fitting of results after \SI{200}{\minute} are not considered artificial, therefore they are not displayed here. 
We can find the linewidth of the methanol close to melting point is almost 2 times larger than that at room-temperature. In other words, $T_2^*$ time decreased by 50\,\%. Besides, the fluctuation of signal frequency is within $\pm \SI{1}{\Hz}$ except the measurement done during the freezing. Considering the spectral resolution here is \SI{1}{\Hz}, and linewidth is 10 times larger than such fluctuation, we can conclude that leading field remained stable during the measurement. 

\section{Noise analysis}\label{sec:Noise_analysis}

Here we characterize noise by the standard deviation $\sigma$ of PSD in the unit of $\Phi_0^{2}/\SI{}{\Hz}$. Notice that $\sigma$ is different from the power of noise, which is the integral of PSD spectrum. 
From \SI{1}{\kHz} to \SI{4}{\MHz}, the noise in SQUID has no dependence on frequency (white noise), which is introduced by the electronics in SQUID (like amplifier and integrator). Typical SQUID noise is approximately $10^{-12}\,\Phi_0^{2}/\SI{}{\Hz}$\cite{squidhandbookintro}. 
To figure out the noise of our SQUID in use and compare it to SPN, a 10-hr SQUID measurement was done with methanol sample and no pulsing. Results are listed in Fig.\,\ref{fig:SpinNoise_Results}. We took chunk size for single PSD as \SI{33}{\ms} so that the spectral resolution (\SI{30}{\Hz}) matches the NMR linewidth (see Fig.\,\ref{fig:PSD_and_Exclusion10hr}(a)) since we try to accumulate SPN in one frequency bin. 
\begin{figure}[htbp]
	\centering
	\includegraphics[width=\textwidth,right]{figures/03_Diagnostic/10hrMeasurement_results_20220919.png}
	\caption{(a) First PSD spectrum of the measurement. (b) Average of all spectra. Due to imperfect correction on the digital RC filter, the PSD curve is not level. (c) Average of spectra with background curvature subtracted. (d) Standard deviation of averaged PSD.}
	\label{fig:SpinNoise_Results}
\end{figure}
The standard deviation of the first spectrum is computed to be $\SI{2.2e-12}{}\,\Phi_0^{2}/\SI{}{\Hz}$. This indicates that we are approaching the limit of SQUID noise. 
LIA adopted digital RC filter as the low-pass filter. Due to imperfect correction, the background is not level. By fitting the background to a polynomial and subtracting PSD with it, we tried to get rid of the effect of the filter. 
By averaging the spectra, the noise goes down as $N^{-1/2}$, where $N$ is the number of spectra. This measurement produced over $10^6$ spectra, and the noise of averaged spectrum is nearly $10^{-15}\,\Phi_0^{2}/\SI{}{\Hz}$. This value is 4 orders larger than what we expected for SPN in \ref{sec:SR:FluxinSQUID:SPN}. Hence the white noise of SQUID is the dominant noise source. 

